\section{Evaluation}
\label{sec:evaluation}

We evaluate identity association on 21 manually curated development videos.  The
benchmark fixes the universe of 41,079 saved, non-interpolated detections and gives
every method the same bounding boxes.  Of these, 41,004 belong to retained surfer
identities and 75 are annotated for rejection.  This deliberately isolates
association from detector recall: the results answer whether the available
observations were grouped correctly, not whether YOLO found every surfer.

The saved observations were selected during manual trajectory reconstruction and
their boxes were Rauch--Tung--Striebel smoothed afterwards.  The benchmark is
therefore conditional on gold-derived observations rather than an end-to-end MOT
test.  It is also an in-sample development-set evaluation: algorithm design was
informed by these videos, and the provenance of the original ILP parameter tuning is
not fully recorded.  We report the results as engineering evidence and failure
analysis, not as a held-out leaderboard.

\subsection{Metrics}

For retained detections, pairwise precision and recall compare every unordered pair.
A pair is positive when both observations have the same annotated identity:

\begin{align}
  P_{\mathrm{pair}} &=
    \frac{\text{correct same-identity pairs}}
         {\text{all pairs joined by the tracker}}, \\
  R_{\mathrm{pair}} &=
    \frac{\text{correct same-identity pairs}}
         {\text{all annotated same-identity pairs}}.
\end{align}

False merges reduce precision and fragmentation reduces recall.  Missing outputs are
treated as singleton predictions rather than being silently omitted.  Micro scores
pool all pairs and consequently weight long identities quadratically; macro scores
average per-video values.  We additionally report contaminated predicted tracks,
which contain more than one gold identity; fragmentation excess, the number of
predicted pieces beyond one per gold identity; and exact videos, whose complete
partition matches the annotation up to track-label permutation.  These absolute
counts matter because a high aggregate F1 can conceal one application-breaking
identity switch.

\subsection{Fixed-observation comparison}

\Cref{tab:association-benchmark} compares the production pipeline with OC-SORT and
BoT-SORT.  OC-SORT receives the fixed detections directly.  BoT-SORT uses its ECC
camera compensation but no pedestrian-specific ReID network.  The production method
reruns conservative preprocessing, masked video stabilization, and offline ILP
association.  Its sail descriptor uses border-derived foreground masking and a
whole-fragment appearance prototype; motion gating evaluates centre, width, and
height; and implausible motion, area, or aspect-ratio transitions are excluded before
optimization.

\begin{table*}[t]
  \centering
  \small
  \caption{Fixed-observation identity-association results on 21 development videos.
  ``Contam.'' counts predicted tracks containing more than one gold identity.}
  \label{tab:association-benchmark}
  \begin{tabular}{lrrrrrrr}
    \toprule
    Method & Coverage & Pair P & Pair R & Pair F1 & Macro F1 & Contam. & Frag. excess \\
    \midrule
    OC-SORT & 0.981 & \textbf{0.985} & 0.662 & 0.792 & 0.857 & 9/165 & 845 \\
    BoT-SORT, no ReID & 0.996 & 0.920 & 0.752 & 0.828 & 0.858 & 16/142 & 244 \\
    Production offline ILP & \textbf{0.999} & 0.957 & \textbf{0.918} & \textbf{0.937} & \textbf{0.970} & 9/95 & \textbf{42} \\
    \bottomrule
  \end{tabular}
\end{table*}

The production system gives the strongest overall association score and reduces
fragmentation by a factor of approximately six relative to BoT-SORT and by a factor
of twenty relative to OC-SORT.  Eleven of 21 videos are exact, compared with
four for BoT-SORT and one for OC-SORT.  OC-SORT attains the highest pairwise precision
by remaining highly conservative, but its 845 excess fragments do not provide the
continuous identity track required by the application.  The production result is a
more useful operating point for this dataset, but its nine contaminated tracks mean
that it still does not satisfy the nominal zero-false-merge requirement.

Runtime is included only as an operational diagnostic.  On the evaluation machine,
OC-SORT took 9.0 seconds, BoT-SORT 497.7 seconds, and the production association path
402.3 seconds.  The workloads are not identical: OC-SORT needs no video decode,
BoT-SORT includes decode and ECC, and production includes masked stabilization but
excludes the initial crop-embedding extraction.

\subsection{Safeguard ablation}

The initial production reconstruction exposed a concentrated false-link problem.
Preprocessor fragments were almost always identity-pure, but the ILP sometimes joined
two pure fragments from different surfers.  Many of the unexpected short-gap errors
occurred at crossings, where the endpoint crops contained both sails.  The previous
appearance model emphasized precisely those boundary frames, the color descriptor's
intended foreground mask was bypassed by an early return, and the ILP requested a
position-only projection from a Kalman model that already represented width and
height.

\Cref{tab:safeguard-ablation} records two staged revisions.  Because several changes
are bundled in the first stage, this is a system ablation rather than a claim that any
single component produced the full difference.  Foreground masking, whole-fragment
appearance aggregation, and four-dimensional motion evaluation raise micro F1 from
0.893 to 0.928.  Adding conservative vetoes for motion NLL and threefold area or
aspect-ratio discontinuities raises precision further and halves contaminated tracks,
without increasing fragmentation.

\begin{table}[t]
  \centering
  \scriptsize
  \caption{Staged production-pipeline revision on the same development set.}
  \label{tab:safeguard-ablation}
  \begin{tabular}{@{}lrrrrr@{}}
    \toprule
    Stage & P & R & F1 & Cont. & Frag. \\
    \midrule
    Original & 0.907 & 0.880 & 0.893 & 20 & 48 \\
    + mask, mean, 4D & 0.939 & 0.917 & 0.928 & 17 & 44 \\
    + conservative vetoes & \textbf{0.957} & \textbf{0.918} & \textbf{0.937} & \textbf{9} & \textbf{42} \\
    \bottomrule
  \end{tabular}
\end{table}

The vetoes are not a second tracker layered on top of the Kalman filter.  The Kalman
state already contains centre, width, and height; the earlier ILP call discarded the
size dimensions.  The revised system uses the full four-dimensional Mahalanobis
evidence and additionally refuses links whose absolute evidence is implausible.  This
second step is necessary because the ILP's start and end penalties can otherwise make
even a very poor link cheaper than splitting two central fragments.

A manual audit found nine wrong selected fragment joins.  Unlike the removed
geometric failures, they concentrate at crossings, visually near-identical sails,
and crops containing both competitors.  Their implications and the unsuccessful
global-margin abstention experiment are discussed in \Cref{sec:discussion}.
